% =============================================================================
%  solver_comparison_tutorial.tex
%  Section of the TUTORIAL paper (paper_tutorial.tex, \input by the lead):
%  "How coupledFoam compares to commercial and other open-source solvers".
%
%  Needs the tutorial preamble (booktabs, longtable, array, tcolorbox boxes
%  intuition / whyitmatters / pitfall, \cfnum, \code, natbib). Measured
%  coupledFoam numbers come only from the generated macros (\cfnum{...});
%  no measured number is typed by hand. Vendor facts carry a citation; what
%  is not publicly documented is marked \cmpnpd.
%  Bibliography keys of this section are prefixed cmp- (references.bib).
% =============================================================================

% --- local helpers (guarded, so the file can be \input more than once) ------
\providecommand{\cmpnpd}{\textit{(n.p.d.)}}
\providecommand{\cmphead}[1]{\multicolumn{6}{l}{\rule{0pt}{2.6ex}\textbf{#1}}\\*[1pt]}
\makeatletter
\@ifundefined{cmpf1box}{%
  \newtcolorbox{cmpf1box}[1]{cfbase, colback=cfwhybg, colframe=cfwhy,
    title={#1}}}{}
\makeatother

% =============================================================================
\section{How coupledFoam compares to commercial and other open-source solvers}
\label{sec:compare}
% =============================================================================

Users of Simcenter STAR-CCM+ or Ansys Fluent have met coupled solvers
before: both products offer one, and so does ENGYS HELYX,
which many motorsport teams use on top of an OpenFOAM-derived code base.
This section puts coupledFoam next to them. For every aspect it says what
coupledFoam does the same, what it does differently, \emph{why}, and when
that should make it better or worse.

\begin{pitfall}{what can and cannot be known about commercial solvers}
  The commercial codes are closed. Their user guides and theory guides
  describe the \emph{method} and the \emph{user settings} (for example the
  default Courant number), but rarely the internals (which smoother runs on
  which level, which operations run in single precision, how a rotating
  zone enters the matrix). Everything below is taken from vendor
  documentation, vendor blog posts or peer-reviewed papers, and each
  statement carries its source. Where a detail is not publicly documented,
  the table says \cmpnpd{} (``not publicly documented''), and the text says
  what typical practice suggests; nothing is filled in by guessing. Vendor
  speed-up claims are the vendors' own measurements, on their own cases,
  against their own baselines. They are quoted as claims, not as verified
  facts.
\end{pitfall}

% -----------------------------------------------------------------------------
\subsection{Four pairs of terms}\label{sec:cmp-words}

\paragraph{Segregated versus coupled.} A \emph{segregated} solver solves
one equation at a time: first momentum with the old pressure, then a
pressure(-correction) equation, then turbulence, and repeats
(Section~\ref{sec:simple}). A \emph{coupled} solver puts momentum and
continuity into one linear system and solves for velocity and pressure
together (Section~\ref{sec:coupled}). In every solver discussed here,
``coupled'' means \emph{velocity and pressure} (and energy, where it
applies). Turbulence is solved separately, after the flow update, in all of
them, commercial ones included.

\paragraph{Pressure-based versus density-based.} These are two historical
families of coupled solvers. A \emph{pressure-based} solver takes pressure
as the unknown that enforces mass conservation; it was born for
incompressible flow and extended to compressible flow. coupledFoam, the
Fluent pressure-based coupled solver~\cite{cmp-ChenPrzekwas2010,cmp-Fluent12TG}
and HELYX-Coupled's pressure-based solver~\cite{cmp-HelyxCoupled} belong
here. A \emph{density-based} solver was born for compressible, high-speed
flow: it solves the conservation laws of mass, momentum and energy as one
vector of equations. At low speed the sound waves are much faster than the
flow and such a solver would crawl, so it uses \emph{low-Mach
preconditioning}, which rescales the pseudo-time behaviour so that sound
and flow travel at comparable speeds~\cite{cmp-WeissSmith1995,Weiss1999}.
Siemens describes the STAR-CCM+ Coupled Flow solver as a density-based
solver with Weiss--Smith preconditioning, used also for constant-density
flow~\cite{cmp-StarCoupledDoc,cmp-StarCoupledGPU2023}.

\begin{intuition}{two ways to the same valley}
  Pressure-based and density-based solvers solve the same physics and,
  when converged on the same mesh with equivalent schemes, should give
  practically the same answer. They differ in the route: how the
  equations are written for the iteration, and what keeps pressure and
  velocity from decoupling. For a steady incompressible car simulation
  both routes are established industrial practice.
\end{intuition}

\paragraph{Pseudo-transient versus relaxed.} A steady answer can be
approached in two ways. \emph{Under-relaxation} (SIMPLE) accepts only a
fixed fraction of each change. \emph{Pseudo-transient continuation}
marches in a fake time towards the steady state, with a time step set by a
Courant (CFL) number (Section~\ref{sec:ptc}). Mathematically, a pseudo-time
term acts like an \emph{implicit} relaxation whose strength adapts to the
local flow; Fluent's documentation calls its pseudo-time method ``a form of
implicit under-relaxation''~\cite{cmp-FluentPseudoTime}. All coupled
solvers here use some form of pseudo-time; several add an explicit
relaxation on top.

\paragraph{Algebraic multigrid (AMG) and its smoother.} All the solvers
solve their large linear systems with algebraic multigrid
(Section~\ref{sec:amg}). For a coupled solver the AMG must handle small
blocks per cell instead of single numbers (``block'' or ``coupled'' AMG),
and the \emph{smoother} (the cheap local iteration on each level) matters a
lot: Fluent's documentation recommends Gauss--Seidel for scalar,
segregated systems and ILU for block-coupled systems~\cite{cmp-Fluent12UG},
which is the same lesson coupledFoam learnt the hard way
(Section~\ref{sec:smoothers}).

% -----------------------------------------------------------------------------
\subsection{The solvers in one paragraph each}\label{sec:cmp-who}

\begin{description}
\item[coupledFoam] (this project). Steady, incompressible RANS; pressure-based,
  fully implicit $u$--$v$--$w$--$p$ coupling with a $4\times4$ block per
  cell, solved in increment form with pseudo-transient continuation, a
  physicality line search, remediation sets and rollback; FGMRES with its
  own block AMG in mixed precision. An add-on to a stock OpenFOAM v2606,
  GPL-licensed, and young: its open problems are listed openly in the
  project files and in Section~\ref{sec:limits}.
\item[Simcenter STAR-CCM+, Coupled Flow.] Density-based coupled solver with
  Weiss--Smith preconditioned Roe flux, implicit pseudo-time marching
  controlled by a Courant number, AMG linear solver; since version 2020.1
  an automatic CFL control is the default for steady runs, optionally with
  a line search on the explicit relaxation~\cite{cmp-StarCoupledDoc,
  cmp-StarAutoCFL2020,cmp-StarEffortless2020}. Siemens names vehicle
  external aerodynamics as a typical application and has shipped a
  GPU-native version since release 2306~\cite{cmp-StarCoupledGPU2023}.
\item[Simcenter STAR-CCM+, Segregated Flow.] The SIMPLE-type alternative
  in the same product, with SIMPLEC available since release 2302, where
  the pressure under-relaxation is 1 and, according to Siemens, the
  velocity under-relaxation can often be raised to 1
  too~\cite{cmp-StarSIMPLEC2023}. It is the common choice for
  incompressible external aerodynamics in STAR-CCM+ (typical practice; the
  vendor does not publish a ``recommended for F1'' statement).
\item[Ansys Fluent, pressure-based coupled.] Solves momentum and the
  pressure-based continuity equation together, with implicit pressure
  gradient and implicit face mass flux including the Rhie--Chow pressure
  dissipation, in $\delta$-form, with a coupled AMG~\cite{cmp-Fluent12TG,
  cmp-ChenPrzekwas2010}. Steady runs are controlled by a ``flow Courant
  number'' (default 200) and explicit relaxation factors for momentum and
  pressure (default 0.75)~\cite{cmp-Fluent12UG}; a separate pseudo-time
  method with automatic global time step is available, and a local
  pseudo-time-step method for this solver appears in recent (beta)
  documentation~\cite{cmp-FluentPseudoTime,cmp-FluentLocalPseudo}. Fluent
  also has a density-based coupled solver, mainly for compressible flow.
\item[ENGYS HELYX-Coupled.] An add-on to the OpenFOAM-derived HELYX
  product: a fully implicit, pressure-based, block-coupled solver for
  incompressible and compressible flow (plus a density-based one),
  local time stepping with a per-cell CFL target, and an optional
  end-to-end single-precision (FP32) mode~\cite{cmp-HelyxCoupled,
  cmp-HelyxTechniques}. ENGYS states that seven of ten F1 teams use HELYX
  and all of them use HELYX-Coupled~\cite{cmp-HelyxCoupled} (vendor claim).
  Linear solver and discretisation internals are not public.
\item[foam-extend \code{pUCoupledFoam}.] The open-source academic
  predecessor closest to coupledFoam: a steady incompressible
  $4\times4$ $\Uv$--$p$ block solver in foam-extend with Rhie--Chow
  interpolation, implicit velocity relaxation (e.g.\ 0.85--0.995, none on
  pressure) and segregated turbulence~\cite{cmp-Jareteg2014,cmp-Jareteg2015};
  its block-selective AMG (classical AMG on a ``primary'' matrix, best with
  the pressure equation) and its parallel version are by Uroi\'c and
  Jasak~\cite{Uroic2018,Uroic2021}. The OpenFOAM coupled solver of Mangani,
  Buchmayr and Darwish~\cite{Mangani2014a,Mangani2014b} is the other
  well-documented line of work.
\item[OpenFOAM \code{simpleFoam}.] The segregated SIMPLE/SIMPLEC baseline
  in current use for the F1 cases~\cite{Patankar1972,VanDoormaal1984,OpenFOAMv2606}, and the
  reference every coupledFoam test is measured against.
\end{description}

% -----------------------------------------------------------------------------
\subsection{The big table}\label{sec:cmp-table}

Table~\ref{tab:cmp} compares the five solvers aspect by aspect.
\cmpnpd{} means ``not publicly documented''. A number without a citation in
the coupledFoam column is a coupledFoam default and can be found in the
sections of this document that are referenced in the text below. After the
table, every group of rows is explained in plain words.

{\scriptsize
\setlength{\tabcolsep}{2.5pt}
\setlength{\LTcapwidth}{\linewidth}
\renewcommand{\arraystretch}{1.15}
\begin{longtable}{>{\raggedright\arraybackslash}p{0.125\linewidth}
                  >{\raggedright\arraybackslash}p{0.165\linewidth}
                  >{\raggedright\arraybackslash}p{0.165\linewidth}
                  >{\raggedright\arraybackslash}p{0.165\linewidth}
                  >{\raggedright\arraybackslash}p{0.155\linewidth}
                  >{\raggedright\arraybackslash}p{0.155\linewidth}}
\caption{coupledFoam next to industrial and open-source solvers for steady
incompressible RANS. \cmpnpd: not publicly documented. Sources: STAR-CCM+
\cite{cmp-StarCoupledDoc,cmp-StarAutoCFL2020,cmp-StarEffortless2020,
cmp-StarCoupledGPU2023,cmp-VolupeCoupled}; Fluent
\cite{cmp-Fluent12UG,cmp-Fluent12TG,cmp-FluentPseudoTime,
cmp-ChenPrzekwas2010}; HELYX \cite{cmp-HelyxCoupled,cmp-HelyxTechniques};
simpleFoam \cite{OpenFOAMv2606}; coupledFoam: this document.}
\label{tab:cmp}\\
\toprule
aspect & coupledFoam & STAR-CCM+ Coupled Flow & Fluent pressure-based coupled
  & HELYX-Coupled (pressure-based) & OpenFOAM simpleFoam\\
\midrule
\endfirsthead
\multicolumn{6}{l}{\small\textit{Table~\ref{tab:cmp}, continued}}\\
\toprule
aspect & coupledFoam & STAR-CCM+ Coupled Flow & Fluent pressure-based coupled
  & HELYX-Coupled (pressure-based) & OpenFOAM simpleFoam\\
\midrule
\endhead
\midrule
\multicolumn{6}{r}{\small\textit{continued on the next page}}\\
\endfoot
\bottomrule
\endlastfoot
%
\cmphead{A. What is solved together}
coupling family
  & pressure-based, fully implicit
  & density-based with low-Mach (Weiss--Smith) preconditioning, also for
    constant density
  & pressure-based, fully implicit
  & pressure-based, fully implicit (a density-based variant also exists)
  & pressure-based, segregated (SIMPLE / SIMPLEC)\\
unknowns in one system
  & $u,v,w,p$: one $4\times4$ block per cell
  & continuity, momentum, energy (and species) as one vector
  & momentum + pressure-based continuity
  & momentum + pressure (``block-coupled'')
  & none: $U$ components, then $p$, one after the other\\
form of the unknown
  & increment $\delta\vect x$ with residual in double
  & pseudo-time increment \cmpnpd{} in detail
  & $\delta$-form (documented)
  & \cmpnpd
  & the new values (with relaxation)\\
pressure--velocity link on the faces
  & Rhie--Chow, implicit $p$--$p$ block; fluxes recomputed with the same
    formula; tensorial variant being added
  & preconditioned Roe flux-difference dissipation (not Rhie--Chow)
  & Rhie--Chow type, implicit in the face mass flux
  & \cmpnpd{} (Rhie--Chow type expected for a pressure-based FV code)
  & Rhie--Chow type momentum interpolation (\code{rAU}, \code{consistent}
    for SIMPLEC)\\
%
\cmphead{B. The steady iteration}
path to steady state
  & pseudo-transient continuation, local time step per cell
  & implicit pseudo-time, Courant number; automatic CFL default for
    steady since 2020.1
  & ``flow Courant number'' (default 200); optional pseudo-time method
    (automatic global step, local step in beta)
  & local time stepping with a per-cell CFL target
  & relaxation only (no pseudo-time)\\
CFL / step control
  & adaptive CFL 1--500 (start 5), grows with residual drop, cut by the
    safety checks; per-cell limiter
  & automatic (algorithm \cmpnpd); expert driver as older alternative
  & fixed user value, or automatic pseudo time step (length/velocity
    scales, factor)
  & automatic per cell (algorithm \cmpnpd)
  & n/a\\
explicit relaxation / line search
  & none on $\Uv,p$; physicality line search $\omega\le1$
    (30\,\% of $U_\text{ref}$)
  & explicit relaxation, constant or adapted by line search
  & explicit relaxation 0.75 default (0.5 for high order, 0.25 for very
    skewed meshes suggested)
  & \cmpnpd
  & typical $p$ 0.3 / $U$ 0.7 (SIMPLE); up to 1 / 0.9 (SIMPLEC)\\
start-up help
  & \code{potentialFoam} start, first-order start-up phase
  & grid-sequencing initialisation
  & hybrid / FMG initialisation (Fluent option; details \cmpnpd)
  & \cmpnpd
  & \code{potentialFoam}, user tricks (upwind first)\\
%
\cmphead{C. The linear solver}
Krylov method
  & FGMRES(30), BiCGStab option
  & \cmpnpd{} (AMG as solver; Krylov acceleration option \cmpnpd)
  & AMG as solver (Krylov acceleration \cmpnpd{} for default)
  & \cmpnpd
  & PCG/PBiCGStab or GAMG as solver, per equation\\
inner tolerance
  & adaptive (Eisenstat--Walker, 0.5 to $10^{-3}$)
  & user convergence tolerance per AMG solve (default \cmpnpd)
  & user termination criterion per AMG solve
  & \cmpnpd
  & fixed \code{relTol}/\code{tolerance} per equation\\
AMG type
  & own block additive-correction AMG, $4\times4$ blocks, pair
    agglomeration on matrix weights
  & AMG (details \cmpnpd)
  & coupled AMG, additive correction (Hutchinson--Raithby), Galerkin
    coarse operator, groups of 2--4 strongest neighbours
  & \cmpnpd
  & scalar GAMG (face-area pair agglomeration) for $p$\\
cycle and smoother
  & V cycle (F, W, K available); damped block ILU(0), relaxation 0.5
  & Gauss--Seidel relaxation documented for AMG in general; coupled
    details \cmpnpd
  & F cycle default for the coupled flow system; ILU smoother default for
    coupled systems (Gauss--Seidel for scalar)
  & \cmpnpd
  & GAMG V-type cycle with Gauss--Seidel or DIC smoother (user choice)\\
precision
  & matrix and Krylov vectors in float; residual, sums, fields in double
  & mixed- and double-precision builds; what runs in single \cmpnpd
  & single- and double-precision modes
  & FP32 end-to-end option, vendor claims about 40\,\% shorter runtimes
  & double (the SPDP build stores fields in float)\\
%
\cmphead{D. Turbulence and robustness}
turbulence
  & segregated: native OpenFOAM model once per outer iteration
  & separate turbulence solvers with own relaxation (typical; coupling
    \cmpnpd)
  & segregated, after the coupled solve (documented)
  & segregated \cmpnpd{} (typical)
  & segregated, relaxed\\
bad cells
  & static set (mesh quality) and dynamic set (behaviour): local upwind,
    smaller local time step, clipping
  & limiters and automatic CFL; per-cell treatment \cmpnpd
  & lower explicit relaxation for skewed meshes; per-cell treatment
    \cmpnpd
  & vendor: ``more permissive with poor quality grids'' than segregated
  & limited / bounded schemes, non-orthogonal limiter, global relaxation\\
failure handling
  & sentinel check, rollback to the last good state, CFL cut, abort
    with \code{\_lastValid} fields
  & \cmpnpd{} (line search on explicit relaxation)
  & \cmpnpd
  & \cmpnpd
  & none (a diverged run is simply lost)\\
oscillation damping
  & optional Anderson acceleration; selective frequency damping being
    added (off by default)
  & \cmpnpd
  & \cmpnpd
  & \cmpnpd
  & none\\
convergence monitoring
  & one \code{CF|} log line per iteration, JSON diagnostics (levels 0--3),
    force-window mean and RMS, stationary-mean stop rule
  & residual and report monitors, stopping criteria (GUI)
  & residual and report monitors, convergence conditions (GUI)
  & residuals, function objects \cmpnpd{} in detail
  & initial residuals, function objects (\code{forceCoeffs},
    \code{solverInfo})\\
%
\cmphead{E. Physics and features}
MRF / rotating wheels
  & MRF, rotation term implicit in the $3\times3$ block
  & supported; implicitness \cmpnpd
  & supported; implicitness \cmpnpd
  & MRF and ``generalised reference frames'' supported
  & MRF, rotation term explicit source\\
sliding and AMI interfaces
  & \code{cyclic}, \code{cyclicAMI} coupled explicitly (lagged)
  & supported \cmpnpd{} in detail
  & supported \cmpnpd{} in detail
  & supported \cmpnpd{} in detail
  & \code{cyclicAMI} implicit inside the linear solver\\
parallel
  & MPI; processor interfaces implicit; coarse-level gathering; tested up
    to 10 ranks so far
  & proven on thousands of cores; GPU-native since 2306
  & proven on thousands of cores
  & vendor: 250\,M cells on about 1000 cores in about 30 minutes
  & proven on thousands of cores\\
transient
  & no (steady only)
  & yes (dual time stepping)
  & yes (coupled or PISO)
  & yes (RANS, URANS, DES, LES)
  & no; \code{pimpleFoam} is the transient sibling\\
compressible
  & no
  & yes, all speeds
  & yes (pressure-based and density-based)
  & yes, up to about Mach 1.5 (pressure-based)
  & no; \code{rhoSimpleFoam} is the sibling\\
%
\cmphead{F. Practicalities}
typical speed
  & measured here against simpleFoam on the test cases, with the same
    criterion for both (Section~\ref{sec:speedtests}); motorbike: see
    Table~\ref{tab:wakespeed}
  & vendor: automatic CFL 2.4$\times$ faster than before on a vehicle case
  & documented: convergence ``significantly improves'' vs segregated;
    memory 1.5--2$\times$
  & vendor: 2--10$\times$ faster than segregated
  & baseline\\
maturity and validation
  & days to weeks old; own test suite T0--T5; open problems listed
  & decades of industrial use and validation
  & decades of industrial use and validation
  & industrial use; vendor names most F1 teams as users
  & very mature, huge user base\\
cost and licence
  & free, GPL-3.0
  & commercial licence
  & commercial licence
  & commercial add-on to an OpenFOAM-derived product
  & free, GPL-3.0\\
transparency
  & full source and every design decision (\code{DECISIONS.md})
  & closed; user guide and theory guide
  & closed; user guide and theory guide
  & closed for the coupled internals
  & full source\\
\end{longtable}
}

% -----------------------------------------------------------------------------
\subsection{Group by group: what is the same, what is different, and why}
\label{sec:cmp-why}

\subsubsection*{A1. Coupling family: pressure-based, like Fluent and HELYX}

\paragraph{Why this choice.} For incompressible flow the density is
constant, so the density-based route needs an artificial device (low-Mach
preconditioning or artificial compressibility~\cite{cmp-WeissSmith1995,
cmp-Economon2020}) to make pressure appear in the time-marching. The
pressure-based coupled formulation works with the incompressible equations
directly; it fits OpenFOAM, whose whole incompressible toolbox (boundary
conditions, turbulence models, MRF, function objects) is written in terms
of $\Uv$, $p$ and the face flux $\phi$. coupledFoam can therefore reuse all
of it unchanged. This is also the choice of Fluent's pressure-based coupled
solver, of HELYX-Coupled and of the academic OpenFOAM coupled solvers
\cite{Darwish2009,Mangani2014a,cmp-Jareteg2014}.

\paragraph{Better or worse, and when.} For steady incompressible car
aerodynamics neither family has a built-in accuracy advantage; the
converged answers depend on the mesh, the schemes and the turbulence model,
not on the solution route. A density-based solver handles compressible
flow naturally, which coupledFoam cannot do at all; for an F1 car at
racing speed (Mach well below 0.3) that does not matter. One practical
difference: STAR-CCM+'s coupled solver stabilises pressure through the
preconditioned Roe dissipation, coupledFoam through the Rhie--Chow term. Both
add a small, mesh-dependent smoothing of the pressure; with the same mesh,
differences should be small but not exactly zero. When the same case is
run in both, forces, not residuals, are the quantities to compare.

\subsubsection*{A2. Pressure--velocity link: Rhie--Chow, like Fluent and
simpleFoam}

\paragraph{Why this choice.} Rhie--Chow~\cite{RhieChow1983} is the
standard for cell-centred pressure-based codes. Fluent documents it
explicitly (``Rhie--Chow pressure dissipation'', implicit in the face mass
flux~\cite{cmp-Fluent12TG}), and simpleFoam's momentum interpolation is the
segregated form of the same idea. coupledFoam puts the direct pressure
difference implicitly into the matrix (Section~\ref{sec:rc}) and uses
\emph{exactly} the same face formula to compute the fluxes after the solve,
so the reported mass imbalance is the continuity residual itself.

\paragraph{Better or worse, and when.} Same idea as the commercial
pressure-based solvers; no advantage either way is expected. A known
weakness of the plain (scalar) Rhie--Chow coefficient is that it ignores
the directional structure of the momentum matrix, which matters in rotating
MRF zones and in stretched boundary-layer cells. A tensorial coefficient
(amendment C2 of the specification) is being added for exactly that reason.
Whether the commercial codes use a scalar or a tensorial coefficient is
\cmpnpd.

\subsubsection*{B1. The steady iteration: pseudo-time instead of
relaxation factors}

\paragraph{What is the same.} Every coupled solver in the table marches in
pseudo-time with a Courant number: STAR-CCM+ (automatic CFL since
2020.1)~\cite{cmp-StarEffortless2020}, Fluent (flow Courant number, optional
automatic pseudo-time step)~\cite{cmp-Fluent12UG,cmp-FluentPseudoTime},
HELYX-Coupled (local time step per cell aiming at an optimal
CFL)~\cite{cmp-HelyxTechniques}. coupledFoam does the same with a local time
step per cell and an adaptive global CFL (Section~\ref{sec:ptc}). The idea
that the CFL should grow as the residual falls is standard in the
literature~\cite{MulderVanLeer1985,Kelley1998,Ceze2013AIAA}.

\paragraph{What is different.} Fluent and STAR-CCM+ \emph{also} apply an
explicit relaxation to the coupled update (Fluent: 0.75 by default,
recommended lower for high-order schemes and skewed
meshes~\cite{cmp-Fluent12UG}; STAR-CCM+: constant or adapted by a line
search~\cite{cmp-StarAutoCFL2020}). coupledFoam has no fixed explicit
relaxation. It uses a \emph{physicality} line search instead: the full step
is taken unless some cell would change by more than a fixed fraction of the
reference velocity or pressure, and only then is the step shortened
(Section~\ref{sec:ls}). A decision record (D-051) explains why a stacked
relaxation was not adopted: the pseudo-time step already plays that role,
and the increment form is needed for the line search, the adaptive inner
tolerance and the rollback.

\paragraph{Why.} A fixed relaxation factor is a global brake: it slows
every cell because a few might misbehave. The line search brakes only when
a step is actually dangerous, and the remediation sets
(Section~\ref{sec:remediation}) brake only the cells that cause it. STAR-CCM+'s
2020 move to an automatic CFL with a line search on the relaxation
factor~\cite{cmp-StarAutoCFL2020} goes in the same direction: less hand
tuning, more automatic step control.

\paragraph{Better or worse, and when.} On clean meshes and well-behaved
flows the automatic step control should give fewer outer iterations than a
conservative fixed setting, without tuning by the user. Where the flow has a
physical tendency to oscillate (bluff-body wakes, massive separation), a
very large pseudo-time step removes the damping that a relaxed solver has,
and the iteration can fall into a \emph{limit cycle}. coupledFoam shows
exactly this on the separated airfoil T3 with $k$-$\omega$ SST, where
\code{simpleFoam} reaches a steady answer and coupledFoam does not
(Section~\ref{sec:limits}). The commercial solvers' answer to such cases is
typically a lower CFL or a lower relaxation set by the user; how their
automatic controllers behave in this situation is \cmpnpd. coupledFoam's
answer is being built: selective frequency damping~\cite{cmp-Akervik2006}
(damps the oscillation and vanishes at the steady state) as an option.

\begin{whyitmatters}{the ``Courant number'' is not the same number everywhere}
  A Fluent flow Courant number of 200, a STAR-CCM+ CFL of 100\,000 in a
  vendor example~\cite{cmp-StarAutoCFL2020} and a coupledFoam CFL of 500
  are not comparable. Each code defines its time step differently (global
  or local, convective or with viscous limits, density-based acoustic
  speed or flow speed), and each adds its own relaxation on top. Compare
  how the \emph{residual and the forces} behave, not the CFL values.
\end{whyitmatters}

\subsubsection*{C1. The linear solver: FGMRES and block AMG}

\paragraph{What is the same.} Everyone uses algebraic multigrid on the
coupled system, and the additive-correction type that coupledFoam uses
(coarse equations are sums of fine equations,
Section~\ref{sec:amg}) is the one Fluent documents for its coupled
AMG~\cite{cmp-Fluent12TG,Hutchinson1986,Weiss1999}. Like Fluent, coupledFoam
keeps the four equations of a cell together when it builds coarse levels,
groups cells with their strongest neighbours (coupledFoam: pair
agglomeration on matrix-based weights, D-039), and smooths with an ILU-type
method: Fluent's documentation names ILU as the default smoother for
coupled systems because it ``has better smoothing properties, especially
for block-coupled systems''~\cite{cmp-Fluent12TG,cmp-Fluent12UG}. The
coupledFoam project reached the same conclusion by measurement: block
Gauss--Seidel failed, block ILU(0) worked (Section~\ref{sec:smoothers}).

\paragraph{What is different.}
\begin{itemize}
\item coupledFoam wraps the AMG in a Krylov method (FGMRES with restart 30)
  and chooses the inner tolerance adaptively
  (Eisenstat--Walker~\cite{EisenstatWalker1996},
  Section~\ref{sec:ew}). Fluent and STAR-CCM+ let the user set a fixed AMG
  termination criterion; whether a Krylov accelerator runs by default in
  their coupled solvers is \cmpnpd.
\item coupledFoam's default cycle is V with a \emph{damped} ILU(0)
  smoother (factor 0.5, D-043). Fluent's default for the coupled flow
  system is the F cycle~\cite{cmp-Fluent12UG}. The V cycle was chosen by a
  measured grid of variants on systems dumped from T1--T4 (D-043); the
  damping was needed because the undamped ILU(0) overshoots on the rows of
  the coupled system that are far from diagonally dominant.
\item The foam-extend line of work uses a different AMG family, a
  \emph{selective} (classical) AMG that coarsens on a scalar ``primary''
  matrix, best taken from the pressure equation~\cite{Uroic2018,Uroic2021}.
  coupledFoam lists it as a Phase-2 preconditioner.
\end{itemize}

\paragraph{Better or worse, and when.} The commercial AMG
implementations have been tuned for decades on millions of customer cases;
coupledFoam's block AMG has been tuned on a handful of test meshes. On the
graded, high-aspect-ratio meshes of T1--T4 its settings had to be
re-tuned several times (D-038, D-039, D-043), and on the finer snappy mesh
T4b the solve still degrades at moderate CFL (Section~\ref{sec:limits}).
The linear solver is expected to be the component where a commercial code
is most clearly ahead in robustness today. Where coupledFoam can be ahead is
cost per linear iteration: hand-written, bandwidth-bound block kernels in
single precision (next item).

\subsubsection*{C2. Precision: float matrix, double residual}

\paragraph{What is the same.} Single precision is industrial practice:
Fluent and STAR-CCM+ both ship single-/mixed-precision and double-precision
executables, and HELYX-Coupled offers an end-to-end FP32 mode that ENGYS
says gives ``around 40\,\% shorter runtimes''~\cite{cmp-HelyxTechniques}.
The reason is the same in all codes: sparse linear algebra waits for
memory, and float halves the data (Section~\ref{sec:precision}).

\paragraph{What is different.} coupledFoam's split is explicit and
documented: the big matrix and the Krylov/multigrid vectors are float;
every residual, sum, norm and reduction, and all fields, are double
(D-001, D-012). Because the solver works in increment form, the final
accuracy is set by the double residual, not by the float solve (classical
mixed-precision iterative refinement~\cite{Baboulin2009,Carson2018}). Which
parts of the commercial codes' single-precision modes are single is
\cmpnpd.

\paragraph{Better or worse, and when.} Pure single precision can stall the
residual at a level set by round-off, which matters when deep convergence
or small force differences between two design variants are required. The
double-residual design avoids that floor at little cost. Against a
commercial double-precision run, coupledFoam should need less memory
bandwidth per linear iteration; against a commercial single-precision run,
it should reach deeper residuals. Neither statement has been measured
against a commercial code in this project.

\subsubsection*{D1. Turbulence: segregated everywhere}

\paragraph{What is the same.} All solvers in the table solve the
turbulence equations \emph{after} the coupled flow update, with their own
relaxation. For Fluent this is documented: the remaining equations are
solved ``in a decoupled fashion as in the segregated
algorithm''~\cite{cmp-Fluent12TG}; the same holds for its density-based
solver. For STAR-CCM+ the turbulence solvers appear as separate solver
objects with their own relaxation (typical practice; the detailed coupling
is \cmpnpd). foam-extend \code{pUCoupledFoam} calls
\code{turbulence->correct()} after the block solve~\cite{cmp-Jareteg2015},
exactly like coupledFoam, which uses the unmodified OpenFOAM models
(Section~\ref{sec:turb}).

\paragraph{Why.} Coupling $k$ and $\omega$ into the block would make it
$6\times6$, more than double the matrix memory, and tie the solver to one
turbulence model. Keeping them segregated lets any OpenFOAM RANS model run
unchanged, including the SST settings of an existing case.

\paragraph{Better or worse, and when.} The price is a lag: the flow step
sees the eddy viscosity of the previous iteration. On attached flows this
costs little. On strongly separated flows the lag is one suspect for
coupledFoam's T3 limit cycle (Section~\ref{sec:limits}); the commercial
codes face the same structure and, in practice, the same advice (lower
CFL or more turbulence relaxation on hard separated cases).

\subsubsection*{D2. Robustness: local treatment and rollback}

\paragraph{What is different.} This is where coupledFoam is most
distinctive. It combines
(i) a static set of bad-quality cells and a dynamic set of misbehaving
cells, both treated locally with first-order convection and a smaller
local time step (Section~\ref{sec:remediation}),
(ii) a physicality line search (Section~\ref{sec:ls}), and
(iii) a sentinel with rollback to the last good state, CFL cut, and, if
all fails, a clean stop that writes the last valid fields and marks the
offending cells (Section~\ref{sec:sentinel}). The commercial codes offer
limiters, automatic CFL, lower relaxation and (STAR-CCM+) a continuity
convergence accelerator~\cite{cmp-VolupeCoupled,cmp-Caraeni2013}; whether
they roll back or treat single cells specially is \cmpnpd.

\paragraph{Why.} Real \code{snappyHexMesh} meshes of a car always contain
a few bad cells. A global brake (lower relaxation or CFL everywhere) would
slow the whole run for the sake of a few hundred cells; local treatment
keeps full speed elsewhere. The rollback turns ``the run blew up at
3~a.m.'' into ``the run stepped back, slowed down, and continued'', or at
worst into a stop with a field that shows where the trouble started.

\paragraph{Better or worse, and when.} The remediation sets are an
accuracy trade: first-order convection in a marked cell is less accurate.
The thresholds were therefore chosen with care (D-044: the dynamic set must
not catch the physically correct suction peaks under a wing; D-047: the
static set is kept near 1\,\% of the cells on snappy meshes). The set
sizes are reported in the log. A commercial code that simply runs the whole mesh with a
lower relaxation may be less efficient but has no such local accuracy loss.

\subsubsection*{D3. Convergence monitoring}

\paragraph{What is different.} coupledFoam logs one line per outer
iteration with every control quantity (CFL, line-search factor, cuts,
residuals, inner tolerance, linear iterations, set sizes, rollbacks), and
optionally JSON diagnostics down to single multigrid levels. With force
coefficients it also reports the window mean and RMS, and it has an
optional stop rule on a stationary mean. The commercial codes offer
flexible residual and report monitors and stopping criteria in their GUI;
they expose far less of the solver's internal state.

\paragraph{Why it matters.} For wake-dominated cars the residuals of
\emph{no} steady solver fall to machine zero, and the forces keep swinging
slightly. The project therefore judges such cases by averaged forces over a
window (D-042), which is also what careful commercial users do. The
built-in window statistics remove a post-processing step. The window must
be the same for the two solvers being compared (D-068): a window that
scales with the length of each run favours the solver that is run for
fewer iterations.

\subsubsection*{E1. Rotating wheels (MRF) and sliding interfaces}

\paragraph{What is different.} coupledFoam puts the MRF rotation term into
the $3\times3$ momentum part of the block, i.e.\ implicitly, whereas
\code{simpleFoam} adds it as an explicit source (Section~\ref{sec:turb}).
The implicit form is the natural one for a coupled solver
\cite{Mangani2014b}; how Fluent and STAR-CCM+ treat it internally is
\cmpnpd. On the other hand, coupledFoam couples \code{cyclic} and
\code{cyclicAMI} patches \emph{explicitly} (lagged by one outer iteration),
where \code{simpleFoam} couples AMI implicitly inside its linear solvers.

\paragraph{Better or worse, and when.} Implicit MRF should help
convergence around fast-spinning wheels; this is only checked by an
informational test so far. Explicit AMI coupling slows convergence across
such interfaces; the current F1 setup (MRF wheels, no AMI) is not affected.
For steady runs with rotating meshes and AMI, \code{simpleFoam} or a
commercial code remains the safer choice until implicit AMI coupling
(Phase~2) exists.

\subsubsection*{E2. Scope: steady, incompressible, RANS only}

\paragraph{The honest difference.} The commercial solvers and HELYX-Coupled
cover transient, compressible, heat transfer, multiphase and scale-resolving
(DES/LES) flows in the same solver family. coupledFoam does exactly one
thing: steady incompressible RANS. Transient or DES runs remain with
\code{pimpleFoam} (or a commercial code). This narrow scope is deliberate:
it keeps the code small enough to be understood, tested and changed in a
short time.

\subsubsection*{F1. Speed, maturity, cost and transparency}

\paragraph{Speed: what has been measured.} Only against \code{simpleFoam},
on the project's test cases, on a shared machine
(Section~\ref{sec:speedtests}). Wall-clock speed-ups of the test runs so
far: about \cfnum{SpeedWallTZeroReOneZeroZero}$\times$ on the cavity at
$Re=100$, \cfnum{SpeedWallTOne}$\times$ on T1 and
\cfnum{SpeedWallTTwo}$\times$ on T2 (a lower bound, because
\code{simpleFoam} did not converge there). On the motorbike, with the
same averaging window for both solvers (D-068), the comparison with the
stored reference gives: \cfnum{WakeSpeedNote} (Table~\ref{tab:wakespeed},
with the conditions of the stored reference run). A time ratio below one
means that coupledFoam is slower. On T3 with SST coupledFoam can fall
into a limit cycle. The finer motorbike T4b is the case closest to a car
in the present campaign (the Ahmed body T5 is deferred, D-063), and its
result is the one to watch. The standard numbers of vendor reports (cost
and memory per million cells, residual reduction, parallel efficiency) are
given in Section~\ref{sec:metrics}.

\paragraph{Speed: what the vendors say.} HELYX: 2--10$\times$ faster than
segregated~\cite{cmp-HelyxCoupled}; STAR-CCM+: automatic CFL 2.4$\times$ faster
than its own previous coupled settings on a vehicle case, and SIMPLEC up
to about 40\,\% faster than SIMPLE in its segregated
solver~\cite{cmp-StarEffortless2020,cmp-StarSIMPLEC2023}; Fluent: coupled
convergence ``significantly improves'' over segregated at 1.5--2 times the
memory~\cite{cmp-Fluent12TG}. These numbers cannot be compared with each
other or with those of this project: different cases, baselines, hardware
and stopping criteria. What they do show is that the industry consensus
agrees with the premise of this project: for steady RANS, a well-built
coupled solver beats a segregated one in wall-clock time, at a memory
price.

\paragraph{Maturity.} This is the biggest difference and it favours the
commercial codes by far. STAR-CCM+ and Fluent have been validated for
decades on industrial cases, including motorsport; HELYX-Coupled is used
in F1 according to its vendor. coupledFoam is a research code with a test
suite of six case families and a list of open problems (the T3 limit
cycle, linear-solve degradation at moderate CFL on the fine snappy mesh,
run-to-run differences on several ranks).

\paragraph{Cost and transparency.} Here coupledFoam and \code{simpleFoam}
win: no licence cost, the full source, and, for coupledFoam, a written
record of every design decision with its evidence. When coupledFoam
behaves unexpectedly, the reason can be traced in the source and the
decision log; with a closed code the user depends on vendor support.

\begin{pitfall}{do not mix solvers in one comparison}
  When an F1 variant run in coupledFoam is compared with a baseline run in
  another code (or even in \code{simpleFoam} with different schemes), part
  of the difference comes from the solver's discretisation and
  stabilisation terms, not from the geometry change. Both variants belong
  in the same solver with the same settings, compared by averaged forces.
\end{pitfall}

% -----------------------------------------------------------------------------
\subsection{Other solvers}\label{sec:cmp-others}

\begin{description}
\item[SU2] (open source) solves incompressible flow with a
  density-based coupled approach (a generalised artificial compressibility
  with preconditioning) on vertex-based unstructured
  grids~\cite{cmp-Economon2020}. Mainly used in aerospace and for adjoint
  design optimisation; not a common choice for race-car aerodynamics.
\item[ICSFoam] (open source, OpenFOAM library) targets implicitly coupled,
  mainly density-based simulation of high-speed flows~\cite{Oliani2023}.
\item[The Mangani--Darwish coupled OpenFOAM solver] covers steady,
  rotating-frame and transient (ALE) incompressible
  flows~\cite{Mangani2014a,Mangani2014b,Mangani2017}; background in the
  textbook of Moukalled, Mangani and Darwish~\cite{cmp-Moukalled2016}.
\end{description}
Cadence (Numeca) Omnis, Code\_Saturne and others are not compared here,
because no solid public information on their coupled steady-state solvers
could be collected within this work.

% -----------------------------------------------------------------------------
\begin{cmpf1box}{What this means for the F1 cases}
\begin{itemize}
\item \textbf{Same family as the tools the teams use.} coupledFoam is a
  pressure-based, block-coupled, pseudo-transient solver, the same
  design family as Fluent's pressure-based coupled solver (which also uses
  a coupled AMG with an ILU smoother) and, as far as its public
  description goes (pressure-based, block-coupled, local pseudo-time
  step), HELYX-Coupled. It is not an exotic method, but a young
  implementation of a standard one.
\item \textbf{What should be better than simpleFoam.} Fewer outer
  iterations and no relaxation factors to tune. Whether this turns into
  shorter wall-clock times is measured case by case with the same
  criterion for both solvers (Section~\ref{sec:speedtests}); on the
  motorbike the measured result is given in Table~\ref{tab:wakespeed}. A
  wall-clock advantage on car-like cases is not established.
\item \textbf{What is not yet as good as a commercial code.} Robustness of
  the linear solver on large, stretched snappy meshes; strongly separated
  flows can limit-cycle (T3); memory is higher than \code{simpleFoam}'s;
  validation is limited to the project's test cases.
\item \textbf{What is gained over a commercial code.} An existing OpenFOAM
  case, with its mesh, boundary conditions, SST model and function
  objects, runs unchanged; the code is free and fully transparent; every
  run reports in its log what the solver is doing.
\item \textbf{Practical recommendation.} coupledFoam is suited to be run on
  the steady MRF half-car cases next to \code{simpleFoam}, with judgement
  by averaged forces over a common window and with the CFL number, the
  line-search factor and the remediation set sizes of the \code{CF|} line
  under observation (Section~\ref{sec:f1}). \code{simpleFoam} (or the
  commercial tool in use) remains the reference until T4b has passed.
\end{itemize}
\end{cmpf1box}
